\documentclass[12pt]{article}
\usepackage{tabularx}
\usepackage{lingmacros}
\usepackage{tree-dvips}
\begin{document}
	
	\section*{Micro-StudioPLC Specification}
	
	Open Standard for Inter Controller communications.

	This specification covers the serial communication exchanges between any type of Controllers
	(embedded devices, computers, SOC/SBC, etc ...).

	Generally, MSPLC Software supports Application Class devices and kRT kernel supports any controller
	with at least 1k of Flash and 100 bytes of Ram/Sram, a timer and a serial communication port.

	The present specificaton is OS/Device agnostic, any device corresponding to the present document can
	be interconnected and accessed through the MSPLS protocol.

	V1.0 of the protocol defines a Kernel and User runtime application on a controller.
	The protocol defines 135 base commands and supports arbitrary frames/extended commands.

	The present Specification provides the tools to standardize peripheral access, for example

	\begin{itemize}
		\item A sensor (temperature, hygrometry, accelerometer, ...)
		\item A peripheral (LCD Screen, HMI, ...)
		\item A Generic MCU (microcontroller with Analog and Digital GPIOs)
		\item An embedded controller (Raspberry, BeagleBoard, Arduino, ...)
		\item An industrial controller (Open source and Arduino bases PLCs)
	\end{itemize}

	The Device manufacturers implements some of the specification commands as needed
	and documents their Usage through the Datasheet/Man pages.

	For example a miniman working device such as temperature sensor may be read through SPI Bus with the following command
	\emph{g0ra}: Reads Channel 0 (temperature) as 16 bit signed integer.

	In this example, the Thermometer CI only supports 1 command and 1 channel. Readback the Analog channel 0 provides
	a direct reading of the temperature as 16-bit signed integer.

	% This work is released under the CC0 1.0 Universal (CC0 1.0) Public Domain Dedication.
	% To the extent possible under law, the author(s) have waived all copyright and related or neighboring rights to this work.
	% For details, see https://creativecommons.org/publicdomain/zero/1.0/
	
	\vspace{2em}
	{\small\textit{This document is released under the CC0 1.0 Universal Public Domain Dedication.\\
  	https://creativecommons.org/publicdomain/zero/1.0/}}

	\subsection*{Definitions}
	
	A \emph{Controller} is a generic term for any device which can be addressed via a serial bus.
	
	It can range from simple microcontrollers such as 8 Bits MCU with I2C/SPI to Modbus devices and custom or Open-Source PLC controllers.

	\subsection*{Command list}

	Kernel Commands:\\

	For complex devices, Kernel commands to control and monitor the state of the controller
	can be implemented.\\

\begin{tabular}{ | m{2.5em} | m{7em}| m{2.5em} | m{17em} | } 
	\hline
	Msg & Xch Spec & Hex & Description \\ 
	\hline
	M1 & skst, stopkst & 0x01 & Stop K-ST (Kernel Space scheduler) \\ 
	\hline
	M2 & rkst, runkst & 0x02 & Run K-ST (Kernel Space scheduler) \\
	\hline
	M4 & sys, systat & 0x04 & Read-back systat then reset min/max \\
	\hline
	M8 & clr, clear & 0x08 & Clear Fault and Alarm Words \\
	\hline
	M16 & sust, stopust & 0x10 & Stop U-ST (User Space scheduler) \\
	\hline
	M32 & rust, runust & 0x20 & Run U-ST (User Space scheduler) \\
	\hline
	M64 & sta, status & 0x40 & Return Status, Alarm and Fault words \\
	\hline
\end{tabular}

	\vspace{2em}
	User Commands:\\

	Generic set of commands to Configure, Read and Write data from/to a controller\\

\begin{tabular}{ | m{2.5em} | m{7em}| m{2.5em} | m{17em} | } 
	\hline
	Msg & Xch Spec & Hex & Description \\ 
	\hline
	M128 & g[n]ci & 0x80 & Set GPIO [n] as digital input \\
	\hline
	M129 & g[n]co & 0x81 & Set GPIO [n] as digital output \\
	\hline
	M130 & g[n]ca & 0x82 & Set GPIO [n] as analog input. Server responds with PWM resolution in bits. \\
	\hline
	M131 & g[n]cp & 0x83 & Set GPIO [n] as PWM output. Server responds with PWM resolution in bits. \\
	\hline
	M132 & & & \\
	\hline
	M133 & & & \\
	\hline
	M134 & g[n]cs+data & 0x86 & Configure special. User defined. \\
	\hline
	M135 & g[n]cr & 0x87 & Read-back GPIO [n] Configuration \\
	\hline
	M136 & g[n]cod & 0x88 & Configure GPIO [n] Open Drain (High Z) \\
	\hline
	M137 & g[n]cpu & 0x89 & Configure GPIO [n] Pull Up \\
	\hline
	M138 & g[n]cpd & 0x8a & Configure GPIO [n] Pull Down \\
	\hline
\end{tabular}

\end{document}

M139 &
M140 &
M141 &
M142 &
M143 &
M144 &
M145 &
M146 &
M147 &
M148 &
M149 &
M150 &
M151 &
M152 &
M153 &
M154 &
M155 &
M156 &
M157 &
M158 &
M159 &
M160 &
M161 &
M162 &
M163 &
M164 &
M165 &
M166 &
M167 &
M168 &
M169 &
M170 &
M171 &
M172 &
M173 &
M174 &
M175 &
M176 &
M177 &
M178 &
M179 &
M180 & 
M181 &
M182 &
M183 &
M184 &
M185 &
M186 &
M187 &
M188 &
M189 &
M190 &
M191 &
M192 &
M193 &
M194 &
M195 &
M196 & 
M197 &
M198 &
M199 &
M200 &
M201 &
M202 &
M203 &
M204 &
M205 &
M206 &
M207 &
M208 &
M209 &
M210 &
M211 &
M212 &
M213 &
M214 &
M215 &
M216 &
M217 &
M218 &
M219 & 
M220 & 
M221 &
M222 &
M223 &
M224 &
M225 &
M226 &
M227 &
M228 &
M229 &
M230 &
M231 &
M232 &
M233 &
M234 &
M235 &
M236 &
M237 &
M238 &
M239 &
M240 &
M241 &
M242 &
M243 &
M244 &
M245 &
M246 &
M247 &
M248 &
M249 &
M250 &
M251 &
M252 &
M253 &
M254 &
M255 &
